\documentclass[a4paper]{article}     
\usepackage{amsfonts}
\usepackage{amsmath}
\usepackage{amssymb}
\usepackage{graphicx}
\usepackage{color}

\newcommand{\Q}{\mathbb{Q}}
\newcommand{\R}{\mathbb{R}}
\newcommand{\llim}{\lim\limits}
\newcommand{\dint}{\displaystyle\int}

\begin{document}

\noindent \large Auteur: Yaman Haj Ahmad \\
\noindent \large Studierichting: Informatica\\
\noindent \large Oefeningenreeks 2B nr. 1.5\\

\medskip

\normalsize

druk de volgende uitdrukking uit in x:\\

$\cos (\operatorname{Bgtan} \left(x\right))$\\

Neem $y = \operatorname{Bgtan} \left(x\right)$ waar $y \in \left]-\dfrac{\pi}{2},\dfrac{\pi}{2}\right[$\\

$\Leftrightarrow \tan \left(y\right) = x$\\

$\Leftrightarrow \dfrac{\sin \left(y\right)}{\cos \left(y\right)} = x$\\

$\Leftrightarrow \sin \left(y\right) = \cos\left(y\right) * x$\\

$\Leftrightarrow  \sin^2\left(y\right) = \cos^2\left(y\right) * x^2 $\\

$\Leftrightarrow 1 - \cos^2 \left(y\right) = \cos^2 \left(y\right) * x^2$\\

$\Leftrightarrow \dfrac{1}{\cos^2\left(y\right)} - 1 = x^2$\\

$\Leftrightarrow \dfrac{1}{\cos^2\left(y\right)} = x^2 + 1$\\

$\Leftrightarrow \dfrac{1}{\cos\left(y\right)} = \sqrt{x^2 + 1}$\\

$\Leftrightarrow \cos\left(y\right) = \dfrac{1}{\sqrt{x^2 + 1}}$\\

$\Leftrightarrow \cos\left(\operatorname{Bgtan}\left(x\right)\right) = \dfrac{1}{\sqrt{x^2 + 1}}$

\begin{thebibliography}{999}
%\bibitem{a} Referentie
\end{thebibliography}
\end{document} 